% Section 9.4, from lectures/L09/appendix/c_exercises.md.

\section{Exercises}
\label{c9:sec:exercises}\label{c9:app:c}

\begin{aside}
Nine, ending with the Cross-check. Do them in order; Exercise~\ref{c9:ex:6} onwards need the target
running.

Where an exercise can be checked mechanically, a \textbf{Check yourself} line says how.

\textbf{Solutions.} The exercises that are not Code are answered in \secref{sol:sec:c9}. The Code
exercises have no solution: each is checked by running it, as its Check yourself line says.
\end{aside}

% ------------------------------------------------------------------------------------------
\exercise[fillaccent2]{Sleeping and spinning}{Recall}
\label{c9:ex:1}

\xpart{a} A task is sleeping and another is spinning on a flag. For each, say what state it is in,
whether it is on the run queue, and what it costs.

\xpart{b} Why can a task in \code{TASK_UNINTERRUPTIBLE} not be killed with signal 9? What is it
usually waiting for, and what does its presence in \code{ps} output tell you?

\xpart{c} Which of the two sleeping states should a driver use, and why essentially always?

\xpart{d} \code{schedule()} is called by the incorrect versions in \secref{c9:app:a2} and by the
correct one in \secref{c9:app:a3}. What else has to happen for it to be a sleep rather than a
yield?

% ------------------------------------------------------------------------------------------
\exercise[fillaccent2]{The lost wakeup}{Recall}
\label{c9:ex:2}

The second version from \secref{c9:app:a2}, with the reader already on the queue the handler wakes:

\begin{ccode}
if (fifo_empty()) {
        set_current_state(TASK_INTERRUPTIBLE);
        schedule();
}
\end{ccode}

\xpart{a} Write out the interleaving that loses the wakeup. Be precise about where the interrupt
lands.

\xpart{b} How wide is the window, in instructions? Given a device interrupting every millisecond,
roughly how long before you hit it?

\xpart{c} Which line of \code{___wait_event}'s expansion closes it? Quote it.

\xpart{d} \code{wait_event_interruptible} takes a \textbf{condition}, not a variable. Give two
separate reasons, one about the lost wakeup and one about what happens after a wake.

\xpart{e} The condition is evaluated many times. What two properties must it therefore have?

\xpart{f} A colleague argues the bug is rare because the window is two instructions. What is wrong
with that reasoning?

% ------------------------------------------------------------------------------------------
\exercise[fillaccent3]{Jiffies}{Hand calculation}
\label{c9:ex:3}

This kernel is \code{HZ=250}.

\xpart{a} How long is a jiffy, in microseconds?

\xpart{b} \code{msecs_to_jiffies(1)}, \code{msecs_to_jiffies(4)}, \code{msecs_to_jiffies(10)}: give
all three. Does it round up or down, and why does that choice matter?

\xpart{c} \code{msleep} sleeps for at least \code{msecs_to_jiffies(n)} jiffies and in practice one
more. Predict the measured duration of \code{msleep(1)}, \code{msleep(4)} and \code{msleep(10)}.

\xpart{d} Two of those three are the same number. Which, and what does that mean for a driver that
calls \code{msleep(1)} in a loop expecting a millisecond each time?

\xpart{e} The same driver is rebuilt on a \code{HZ=1000} kernel. Recompute all three.

\xpart{f} \code{jiffies} is compared with \code{time_after} rather than \code{>}. What goes wrong
with \code{>}, and how often, on a 32-bit counter at \code{HZ=250}?

% ------------------------------------------------------------------------------------------
\exercise[fillaccent3]{Choosing a delay}{Hand calculation}
\label{c9:ex:4}

For each, choose \code{udelay}, \code{usleep_range}, \code{msleep}, or a timer, and justify it in
one sentence.

\xpart{a} A hardware register needs 2 us to settle, inside a spinlock.

\xpart{b} Polling a status bit that takes about 500 us, in \code{probe}.

\xpart{c} Retrying a failed transfer in 100 ms.

\xpart{d} Waiting 5 ms in a hard interrupt handler.

\xpart{e} Giving up on a device that has not responded within 2 seconds.

\xpart{f} Waiting 50 us in process context, as accurately as possible.

Then: \textbf{g)} one of the six is a trick. Which, and what is the right answer?

% ------------------------------------------------------------------------------------------
\exercise[fillaccent4]{The four cases of \code{read}}{Design}
\label{c9:ex:5}

Your driver's \code{read} faces four situations. For each, give the return value and say what the
calling program observes.

\xpart{a} 40 samples available, 100 requested.

\xpart{b} Nothing available, descriptor is blocking.

\xpart{c} Nothing available, descriptor has \code{O_NONBLOCK}.

\xpart{d} Nothing available, blocking, and the user presses Ctrl-C.

Then:

\xpart{e} For \textbf{b)}, Chapter~\ref{c5:ch} returned \code{-EAGAIN} and this chapter blocks
instead. What changed, and why was \code{-EAGAIN} the right answer then?

\xpart{f} For \textbf{d)}, what happens if you return \code{-EINTR} instead of \code{-ERESTARTSYS}?
What if you ignore the return value of \code{wait_event_interruptible} entirely?

\xpart{g} What does returning 0 mean, and name a program that would silently do the wrong thing.

% ------------------------------------------------------------------------------------------
\exercise{Block until there is data}{Code}
\label{c9:ex:6}

Write \file{lectures/L09/lab/qa_wait.c}. Start from your Chapter~\ref{c8:ch} driver and add a
character device.
\begin{itemize}
  \item Keep the interrupt handler from Chapter~\ref{c8:ch}, but have it copy samples into a
    driver-side FIFO of 256 entries under a spinlock, and call \code{wake_up_interruptible} when it
    added anything.
  \item Add \file{/dev/qa_wait}, registered the long way as in Chapter~\ref{c5:ch}.
  \item \code{read} returns whole samples. Empty and blocking: wait. Empty and \code{O_NONBLOCK}:
    \code{-EAGAIN}. Interrupted: \code{-ERESTARTSYS}.
  \item Implement \code{poll}.
  \item Take a \code{period_ns} parameter; \textbf{zero means create the device but leave it
    stopped}, which is how the tests reach the idle-device cases.
\end{itemize}

\checkyourself \code{make test L=L09} reports \textbf{PASSED}, with eighteen checks in total: ten
from the harness and eight from the userspace tester across its two modes. The shipped
\code{qa_wait_test} exercises both a running and a stopped device.

\xpart{a} You will be tempted to hold the spinlock across \code{wait_event_interruptible}. What
happens, and which config option this course enables tells you?

\xpart{b} Name your wait queue \code{readq} and try to build. Quote the error and explain it.

\xpart{c} With the device running, drain the device and time the next \code{read}. How long did it
block, and what woke it?

% ------------------------------------------------------------------------------------------
\exercise{Measure the sleeps}{Code}
\label{c9:ex:7}

Write \file{lectures/L09/lab/qa_timing.c}. It must print \code{HZ} and the jiffy length, then
measure the average duration of \code{msleep(1)}, \code{msleep(4)}, \code{msleep(10)},
\code{usleep_range(1000, 1000)}, \code{usleep_range(100, 100)} and \code{udelay(100)} over a
configurable number of repetitions, using \code{ktime_get}.

\checkyourself \code{msleep(1)} and \code{msleep(4)} come out within a few hundred microseconds of
each other, and \code{usleep_range(1000)} comes out below one jiffy.

\xpart{a} Why \code{ktime_get} and not \code{ktime_get_real}?

\xpart{b} \code{udelay(100)} measures about 112 us and \code{usleep_range(100)} about 307 us.
Explain both the ordering and the gap.

\xpart{c} Your module measures \code{msleep} by calling it 200 times in a loop. Name one thing that
measures which a single call would not, and one thing it hides.

% ------------------------------------------------------------------------------------------
\exercise[fillaccent4]{A timer instead of an interrupt}{Design}
\label{c9:ex:8}

Your driver gets its data because the device interrupts. Suppose it did not.

\xpart{a} Rewrite the design using a \code{timer_list} that polls the device every 4 ms. What is the
callback allowed to do, and what would you have to move?

\xpart{b} With \code{hrtimer} instead, what changes about resolution and about where the callback
runs?

\xpart{c} For each design, say what happens to a sample that arrives just after a poll.

\xpart{d} The device's FIFO holds 16 samples. At a 1 ms sample rate, what is the slowest poll
interval that cannot lose data? What does that number become if the FIFO were 4 deep?

\xpart{e} State the general rule for when polling is the right answer and when an interrupt is.

% ------------------------------------------------------------------------------------------
\exercise[fillaccent]{A millisecond that is not a millisecond}{Cross-check}
\label{c9:ex:9}

Predict a duration from \code{HZ}, measure it, and find that one call takes eight times as long as
it asks for.

\xpart{a} \textbf{Predict.} This kernel is \code{HZ=250}. Using \code{msecs_to_jiffies} rounding up,
and the rule that \code{schedule_timeout} guarantees \emph{at least} the requested number of
jiffies, predict the measured duration of:
\begin{itemize}
  \item \code{msleep(1)}
  \item \code{msleep(4)}
  \item \code{msleep(10)}
  \item \code{usleep_range(1000, 1000)}
\end{itemize}
Write all four down before running anything.

\xpart{b} \textbf{Measure.} Load \code{qa_timing} with \code{reps=200} and record all four.

\xpart{c} \textbf{Reconcile.} Answer each:
\begin{itemize}
  \item Which call came out furthest from what it asked for, and by what factor? Did your
    prediction expect it?
  \item Two of the four measurements are the same number to within a rounding error. Which two, and
    is that a bug in the measurement, in \code{msleep}, or in your expectation?
  \item \code{msleep(4)} asks for exactly one jiffy and takes two. Where does the second jiffy come
    from? \code{msleep}'s source does not add one; find the sentence in \code{schedule_timeout}'s
    documentation that explains it, and say why a guarantee of ``at least'' costs an extra tick.
  \item \code{usleep_range(1000)} asks for the same duration as \code{msleep(1)} and takes a fifth
    as long. What is it using instead, and what did you give up to get that?
\end{itemize}

\xpart{d} \textbf{Make a prediction and test it.} Your account of the extra jiffy implies something
specific about \code{msleep(5)} on this kernel. Predict it, measure it, and say whether you were
right.

\xpart{e} \textbf{Change the constant.} Recompute all four predictions for \code{HZ=1000}. Which of
the four measurements would change, and which would not? What does that tell you about which of
these functions are quantised and which are not?

\xpart{f} A driver contains \code{for (i = 0; i < 1000; i++) msleep(1);} in a loop, with a comment
saying ``wait one second''. How long does it actually take on this kernel? On a \code{HZ=1000}
kernel? Rewrite the line so it does what the comment says on both.

\xpart{g} You are asked to document your driver's worst-case response time. Everything measured here
was an average over 200 repetitions on an idle machine. Say what you would put in the document, what
you would refuse to claim from these numbers, and what you would have to measure instead.
Chapter~\ref{c12:ch} is that measurement.
